% BOOK-FILE-ID: f-46762ad20d62
\subsection{差分で線形漸化式を解く}\label{節：差分で線形漸化式を解く}
\sectionref*{節：隣接4項型}の例題(1)では,
\[
  a_{n+3}-3a_{n+2}+3a_{n+1}-a_n=0
\]
を「階差の階差のそのまた階差が0」と読んで解いた.
係数$1,-3,3,-1$は,どこかで見覚えのある数の並びではないだろうか.
ここでは,この考え方がどこまで広げられるかを調べよう.

\paragraphbf{差分を繰り返す}
以下,添字を1つ進める操作を$E$と書く.つまり$Ea_n=a_{n+1}$である.
すると差分は$\Delta a_n=a_{n+1}-a_n$なので,\,$\Delta=E-1$と書ける.
差分を2回,3回と取ると,
\[
  \begin{aligned}
    \Delta^2a_n&=a_{n+2}-2a_{n+1}+a_n,\\
    \Delta^3a_n&=a_{n+3}-3a_{n+2}+3a_{n+1}-a_n
  \end{aligned}
\]
となる.\,$(E-1)^2,\ (E-1)^3$を二項定理で展開したものと同じ係数が並ぶのだ.
先ほどの係数の並びは,符号付きの二項係数である.
なお,\,$\Delta^2a_n$は差分を2回取ったものであり,\,$(\Delta a_n)^2$とは全く別のものである.

では,\,$\Delta^ra_n=0$を満たす数列はどんな数列だろうか.
$\Delta^2a_n=0$なら,\,$\Delta a_n$が一定なので$a_n$は等差数列,つまり$n$の1次以下の式である.
$\Delta^3a_n=0$なら,\,$\Delta^2a_n$が一定,\,$\Delta a_n$が1次以下,\,$a_n$が2次以下と,和を取るたびに次数が1つずつ上がる.
逆に,多項式の差分を取ると次数が1つずつ下がり,定数の差分は0になる.
こうして
\[
  \Delta^ra_n=0
  \iff
  a_n\text{は}n\text{の}r-1\text{次以下の多項式}
\]
が分かる.
和を取るたびに初項が1つずつ必要になるので,\,$a_1,\ldots,a_r$の$r$個の初期値で数列が決まる.

ここで\sectionref*{節：二項係数で多項式を表す}の$\Delta\binom nk=\binom n{k-1}$が役に立つ.
部品$\binom{n-1}{k}$は,差分を取るたびに$\binom{n-1}{k-1}$へと1つずつ前の部品に移るので,
\BookMathLayout{\[
  a_n=a_1+\Delta a_1\binom{n-1}1+\Delta^2a_1\binom{n-1}2+\cdots
\]}{\[
  \begin{aligned}
    a_n&=a_1+\Delta a_1\binom{n-1}1\\
    &\quad+\Delta^2a_1\binom{n-1}2+\cdots
  \end{aligned}
\]}
と,初期値から作る差分をそのまま係数にして一般項が書けるのだ.\footnote{右辺は$\Delta^{r-1}a_1\binom{n-1}{r-1}$までで終わる.
  両辺の差分をそれぞれ取ると同じ形が現れることから確かめられる.}
例えば隣接4項型の例題(1)の$1,2,4$では,\,$\Delta a_1=1,\ \Delta^2a_1=1$なので
\[
  a_n=1+(n-1)+\frac{(n-1)(n-2)}2=\frac{n^2-n+2}2
\]
となり,本編の答えと一致する.
行列の積が簡単になる向きを選んだように,差分が簡単になる部品を選んだのである.

ただし,\,$1,2,4$の続きは漸化式を決めて初めて定まることに注意しよう.
同じ$1,2,4$でも,\,$a_{n+1}=2a_n$なら次は8,\,$\Delta^3a_n=0$なら次は7である.
初めの数項だけから一般項を断定せず,漸化式と初期値から確かめる必要がある.

\paragraphbf{差分で項数を1つ減らせるとき}
隣接4項型の例題(2)では,\,$d_n=a_{n+2}-a_{n+1}-2a_n$と置くと$d_{n+1}-d_n=0$となり,漸化式が1段落ちた.
なぜこの$d_n$で差分がうまく取れたのだろうか.
元の漸化式$a_{n+3}-2a_{n+2}-a_{n+1}+2a_n=0$の特性方程式の左辺を$P(x)$とすると,
\[
  P(x)=x^3-2x^2-x+2=(x-1)(x^2-x-2)
\]
と$x-1$で割り切れる.これを$E$で書けば
\[
  (E-1)(E^2-E-2)a_n=\Delta d_n=0
\]
であり,\,$d_n=(E^2-E-2)a_n$は一定となる.
一般に,定数係数の付加項のない線形漸化式では,
\[
  P(1)=0
\]
すなわち\textbf{係数の和が0}なら,\,$P(x)=(x-1)Q(x)$と書けて,\,$Q(E)a_n$が一定になる.
例題(3)も$1-5+7-3=0$なので,同じように1段落とせる.
逆に,\,$a_{n+2}-5a_{n+1}+6a_n=0$のように係数の和が0でなければ,差分を取っても項数は減らない.
このように差分を因数として取り出し,項数を1つ減らせるのは,特性方程式が1を解に持つ場合である.
ただし,差分そのものはほかの場合にも取ることができる.例えば次に見るように,累乗で割った数列に対して差分を使うこともできる.

\sectionref*{節：特性方程式型}の別解はこの逆向きにあたる.
$a_{n+1}=pa_n+q$の添字をずらして引いた$a_{n+2}-(p+1)a_{n+1}+pa_n=0$は,係数の和が$1-(p+1)+p=0$となる.
定数を消すために差分を取ると,特性方程式の解に1が1つ加わるのである.

\paragraphbf{累乗で割って差分を取る}
特性方程式の解が1でないときはどうだろうか.
本編の隣接3項型で作った$a_{n+1}-\alpha a_n$に注目しよう.
$\alpha\ne0$として,\,$u_n=\frac{a_n}{\alpha^{n-1}}$と置いて$\alpha^{n-1}$を取り除くと,
\[
  \begin{aligned}
    a_{n+1}-\alpha a_n
    &=\alpha^nu_{n+1}-\alpha\cdot\alpha^{n-1}u_n\\
    &=\alpha^n\Delta u_n
  \end{aligned}
\]
となる.\,$\alpha=1$なら,ただの差分である.
つまり,本編で$a_{n+1}-\alpha a_n$を作る操作は,\textbf{$\alpha^{n-1}$で割った数列の差分を取る}ことに対応する.
特性方程式の因数$x-\alpha$の1つが,この操作1回に対応している.

同じ操作をもう一度施すと,\,$\alpha^n\Delta u_n$を$\alpha^{n-1}$で割った数列は$\alpha\Delta u_n$なので,
\[
  (E-\alpha)^2a_n=\alpha^{n+1}\Delta^2u_n
\]
となる.繰り返せば,\,$(E-\alpha)^ra_n=\alpha^{n+r-1}\Delta^ru_n$である.
したがって,特性方程式の左辺が$(x-\alpha)^r$なら,\,$\Delta^ru_n=0$より$u_n$は$r-1$次以下の多項式になり,
\[
  a_n=\alpha^{n-1}\times(r-1\text{次以下の多項式})
\]
を得る.\footnote{$\alpha=0$のときは$\alpha^{n-1}$で割れないので別に扱う.例えば$a_{n+2}=0$は,\,$a_1,a_2$を自由に決めて,その後はすべて0になる.特性方程式の解0は,初めのいくつかの項にしか影響しない.}
本編の重解で$n$が現れたのは,\,$r=2$の場合である.
特性方程式が異なる解を持つときは,それぞれの解からこうして得た解を,\sectionref*{節：答えの形を見通す}の重ね合わせで組み合わせればよい.\footnote{こうして得た解がちょうど初期値の個数だけあり,互いに独立であることの証明には線形代数を使う.}

\paragraphbf{付加項がある場合}
\sectionref*{節：付加項の次数が上がる理由}で見た次数の補正も,同じ計算で説明できる.
付加項$P(n)\alpha^n$の0でない$\alpha$が,付加項のない式の特性方程式の解と$m$重に重なっているとしよう.
$Q(n)\alpha^n$の形の解を探すと,左辺の$(E-\alpha)^m$の部分が$\Delta^m$として$Q$に働き,次数を$m$だけ下げる.
そのため,\,$Q$の次数は$P$より$m$だけ高くなければならない.
例えば$a_{n+2}-4a_{n+1}+4a_n=2^n$では,\,$2$が重解なので,\,$n^2\cdot2^n$の形で探すことになる.

\paragraphbf{2つの解法}
定数係数の線形漸化式には,2つの解法があることが分かった.
\begin{itemize}
  \item \textbf{等比型に分ける}：特性方程式の異なる解ごとに等比型を取り出す.
  \item \textbf{累乗で割って差分を取る}：0でない解の累乗で割り,対応する因数を差分に移す.解が重なるときは,差分を繰り返す.
\end{itemize}
どちらの場合も,特性方程式の解とその重なり方が,一般項の形を教えてくれるのである.
